\Askhsh{33754}{4}
\begin{ylh}{2}
    \item 3.1. Εξισώσεις 1ου Βαθμού
    \item 6.1. Η Έννοια της Συνάρτησης
    \item 6.2. Γραφική Παράσταση Συνάρτησης
\end{ylh}
Για την ενοικίαση ενός συγκεκριμένου τύπου αυτοκινήτου για μια ημέρα, η εταιρία Α χρεώνει τους πελάτες της σύμφωνα με τον τύπο:
\begin{equation*}y = 60 + 0,20x\end{equation*}

Όπου $x$ είναι η απόσταση που διανύθηκε σε km και $y$ το ποσό της χρέωσης σε ευρώ.

\begin{alist}
\item Τι ποσό θα πληρώσει ένας πελάτης της εταιρείας Α ο οποίος, σε μια ημέρα, ταξίδεψε $400Km$;
\monades{5}
\item Πόσα χιλιόμετρα ταξίδεψε ένας πελάτης ο οποίος για μια ημέρα πλήρωσε $150$ ευρώ;
\monades{5}
\item Μια άλλη εταιρεία, η Β, χρεώνει τους πελάτες της ανά ημέρα σύμφωνα με τον τύπο:

\begin{equation*}y = 80 + 0,10x\end{equation*}

όπου, όπως και προηγουμένως, $x$ είναι η απόσταση που διανύθηκε σε km και $y$ είναι το ποσό της χρέωσης σε ευρώ. Να εξετάσετε ποια από τις δύο εταιρείες μας συμφέρει να επιλέξουμε, ανάλογα με την απόσταση που σκοπεύουμε να διανύσουμε.
\monades{10}
\item Αν $f(x) = 60 + 0,20x$ και $g(x) = 0,80 + 0,10x$ είναι οι συναρτήσεις που εκφράζουν τον τρόπο χρέωσης των εταιρειών Α κα Β αντίστοιχα, να βρείτε τις συντεταγμένες του σημείου τομής των γραφικών παραστάσεων των συναρτήσεων $f$ και $g$ και να εξηγήσετε τι εκφράζει η τιμή κάθε μιας από τις συντεταγμένες σε σχέση με το πρόβλημα το ερωτήματος γ).

\monades{5}
\end{alist}

